\section{Анализ предметной области}

\subsection{Особенности предметной области и назначение разработки}
Предметной областью данной работы является разработка игровых движков — программных систем, обеспечивающих создание и функционирование видеоигр и интерактивных приложений в реальном времени. Ключевой особенностью предметной области выступает высокая сложность системной интеграции, поскольку игровой движок должен объединять разнородные подсистемы, такие как рендеринг, физика, звук, обработка ввода и управление ресурсами, каждая из которых предъявляет собственные требования к производительности и синхронизации. Другой важнейшей особенностью являются жёсткие требования к работе в реальном времени: современные игры должны обеспечивать частоту смены кадров не менее 60 кадров в секунду, что отводит на обработку одного кадра не более 16,6 миллисекунды, включая все вычисления физики, логики и графики.

\subsection{Современные архитектурные подходы}
В области разработки игровых движков за последние годы сформировалось несколько архитектурных подходов, каждый из которых имеет свои сильные и слабые стороны. Наиболее традиционным является объектно-ориентированный подход, при котором игровые сущности представляют собой объекты классов, организованные в иерархию наследования. В такой модели базовый класс игрового объекта содержит методы обновления и отрисовки, а конкретные типы объектов (игрок, враг, снаряд) наследуют и переопределяют эти методы. Данный подход интуитивно понятен и широко распространён, однако он приводит к проблеме глубоких иерархий, когда изменение поведения на одном уровне наследования может непредсказуемо повлиять на все дочерние классы. Кроме того, объектно-ориентированный подход затрудняет гибкое комбинирование свойств объекта, что в игровой разработке требуется постоянно, например, когда объект должен одновременно обладать способностью двигаться, атаковать и издавать звук.

В ответ на ограничения объектно-ориентированного подхода получила распространение компонентная архитектура. В этой модели игровой объект представляет собой контейнер, в который можно добавлять отдельные компоненты, отвечающие за конкретные аспекты поведения. Компонент перемещения отвечает за позицию и скорость, компонент рендеринга определяет визуальное представление объекта, компонент здоровья управляет очками прочности, а компонент искусственного интеллекта задаёт логику поведения. Такой подход позволяет гибко конфигурировать объекты без создания сложных иерархий наследования, просто добавляя или удаляя компоненты. Именно компонентную архитектуру используют такие популярные движки, как Unity и Unreal Engine, что подтверждает её эффективность для широкого круга игровых проектов.

\subsection{Обзор графических API для разработки игровых движков}
Для разработки игровых движков, особенно трёхмерной графики, критически важным является выбор графического прикладного программного интерфейса, поскольку именно через этот интерфейс движок взаимодействует с видеокартой и графическим драйвером для отрисовки изображения. Графический API предоставляет разработчику набор функций для работы с буферами вершин, текстурами, шейдерами и другими ресурсами графического конвейера. На сегодняшний день на рынке существуют четыре основных графических API: OpenGL, Vulkan, DirectX 11, DirectX 12 и Metal, каждый из которых имеет свою историю, область применения и архитектурные особенности.

OpenGL является одним из старейших и наиболее распространённых графических API, разработка которого началась в 1992 году компанией Silicon Graphics и впоследствии перешла под контроль некоммерческого консорциума Khronos Group. Основным преимуществом OpenGL является его кроссплатформенность — он поддерживается на операционных системах Windows, Linux, macOS в устаревшем виде, а также на Android и iOS. OpenGL предоставляет разработчику достаточно высокий уровень абстракции, беря на себя управление многими низкоуровневыми аспектами работы графического процессора, такими как управление памятью буферов и синхронизация команд. Это делает OpenGL относительно простым в изучении и использовании по сравнению с более современными API, а порог входа для создания простого рендеринга составляет около двухсот-трёхсот строк кода. Однако эта простота является обратной стороной ограниченного контроля над GPU: разработчик не может явно управлять тем, как драйвер планирует выполнение команд, что в определённых сценариях приводит к неоптимальной производительности, особенно при большом количестве отдельных вызовов отрисовки. Кроме того, OpenGL долгое время страдал от проблемы глобального состояния, когда вызов любой функции мог зависеть от ранее установленных параметров, что усложняло многопоточную работу.

DirectX 11 и DirectX 12 являются графическими API, разработанными корпорацией Microsoft исключительно для операционной системы Windows и игровой консоли Xbox. DirectX 11, выпущенный в 2009 году, по своей концепции близок к OpenGL — он также предоставляет достаточно высокий уровень абстракции и относительно прост в использовании. DirectX 11 широко применяется в коммерческих игровых проектах, особенно тех, которые ориентированы исключительно на платформу Windows, и его производительность на этой платформе зачастую превосходит OpenGL благодаря более плотной интеграции с драйверами Windows. DirectX 12, выпущенный в 2015 году, представляет собой кардинально переработанный API низкого уровня, который предоставляет разработчику гораздо более прямой контроль над графическим процессором. В DirectX 12 разработчик самостоятельно управляет выделением памяти для буферов, синхронизацией между командными очередями, а также распределением работы между центральным и графическим процессорами. Это позволяет достичь значительно более высокой производительности в сложных сценах с тысячами объектов, но цена такого контроля — лавинообразный рост сложности разработки. Для создания минимального работающего приложения на DirectX 12 требуется более тысячи строк кода только для инициализации, причём большая часть этого кода посвящена управлению ресурсами и синхронизацией, а не собственно отрисовке.

Vulkan, также разрабатываемый консорциумом Khronos Group, был выпущен в 2016 году как духовный наследник OpenGL, но с архитектурой, вдохновлённой DirectX 12 и улучшениями, призванными сделать его более кроссплатформенным и эффективным. Как и DirectX 12, Vulkan является низкоуровневым API, предоставляющим разработчику явный контроль над графическим конвейером и памятью GPU. Vulkan поддерживается на Windows, Linux, Android и, через сторонние реализации, на macOS. Ключевая идея Vulkan заключается в том, что многие проверки и операции, которые в OpenGL выполнялись драйвером во время работы программы, в Vulkan переносятся на этап инициализации, что снижает накладные расходы во время рендеринга. Кроме того, Vulkan с самого начала проектировался для многопоточности: разные потоки могут параллельно строить командные буферы, которые затем отправляются на выполнение. Основным недостатком Vulkan является его чрезвычайная сложность — объём шаблонного кода для инициализации может достигать нескольких тысяч строк, а любая ошибка в управлении синхронизацией приводит к трудноотлаживаемым проблемам, вплоть до краша драйвера.

Metal является проприетарным графическим API, разработанным компанией Apple и представленным в 2014 году. Metal поддерживается только на операционных системах Apple: macOS, iOS, tvOS и iPadOS. По своей концепции Metal находится между высокоуровневыми API и низкоуровневыми, предоставляя достаточно прямой контроль над GPU, но с меньшим объёмом шаблонного кода по сравнению с Vulkan или DirectX 12. Для разработчиков, ориентированных исключительно на экосистему Apple, Metal является оптимальным выбором, поскольку он тесно интегрирован с графическими драйверами Apple и позволяет добиться наилучшей производительности на этой платформе. Однако для кроссплатформенных проектов Metal неприменим, поскольку он не работает за пределами устройств Apple.

Сравнивая перечисленные API между собой, можно выделить несколько ключевых критериев выбора. По кроссплатформенности лидируют OpenGL и Vulkan, которые поддерживаются на большинстве настольных и мобильных платформ, тогда как DirectX привязан к Windows и Xbox, а Metal — к экосистеме Apple. По сложности освоения и объёму необходимого кода OpenGL является самым доступным, за ним следуют DirectX 11 и Metal, и значительно более сложными являются DirectX 12 и Vulkan. По уровню контроля над GPU и потенциальной производительности OpenGL уступает остальным, а Vulkan и DirectX 12 предоставляют максимальный контроль. По зрелости и распространённости в игровой индустрии OpenGL и DirectX 11 остаются наиболее укоренёнными, хотя новые проекты всё чаще выбирают Vulkan или DirectX 12 для достижения максимальной производительности.

\subsection{Технологический стек его обоснование}
Для реализации игрового движка в рамках данной дипломной работы был выбран следующий технологический стек. В качестве основного языка программирования используется C\#. Этот язык был выбран по нескольким причинам. Во-первых, C\# предоставляет управляемую среду выполнения с автоматической сборкой мусора, что существенно снижает риск ошибок, связанных с утечками памяти и висячими указателями, по сравнению с языком C++. Во-вторых, C\# обладает богатой стандартной библиотекой, включая удобные коллекции, LINQ для работы с данными, механизмы рефлексии и асинхронного программирования, что ускоряет разработку по сравнению с C++. В-третьих, C\# является кроссплатформенным языком благодаря экосистеме .NET, что позволяет при необходимости собирать движок под различные операционные системы. В-четвёртых, выбор C\# имеет и образовательную ценность, поскольку большинство коммерческих игровых движков, включая Unreal Engine и Godot до версии 4.0, традиционно пишутся на C++, и разработка движка на C\#p позволяет провести сравнительный анализ двух подходов и выявить преимущества и недостатки управляемой среды в контексте графики реального времени.

В качестве графического API выбран OpenGL версии 4.6. Причинами этого выбора являются кроссплатформенность OpenGL, поддерживающего Windows, Linux и macOS, относительно низкий порог входа по сравнению с Vulkan или DirectX 12, а также достаточный уровень контроля над графическим конвейером для реализации всех необходимых функций игрового движка. OpenGL 4.6 включает современные возможности, такие как прямой доступ к состояниям, шейдерные хранимые буферы и вычислительные шейдеры, что позволяет использовать актуальные техники рендеринга. Отказ от Vulkan и DirectX 12 обусловлен их чрезвычайной сложностью, которая привела бы к неоправданному росту объёма шаблонного кода и смещению фокуса с архитектуры игрового движка на низкоуровневое управление ресурсами GPU.

\subsection{Обзор существующих решений}
\subsubsection{Анализ возможностей Unity}
Unity является одним из наиболее распространённых игровых движков в мире, особенно в сегменте мобильной разработки, где, по статистике, около семи из каждых десяти мобильных игр созданы с использованием именно этого инструмента . Такая популярность обусловлена выверенным балансом между функциональностью, простотой освоения и кроссплатформенностью. В рамках анализа предметной области целесообразно рассмотреть ключевые технологические возможности Unity, которые делают его референтным решением для сравнения с разрабатываемым движком, а также выделить его ограничения, требующие понимания при выборе архитектурных решений.

Главным преимуществом Unity является его доступность для широкого круга разработчиков, включая независимых студий и одиночных энтузиастов. Движок использует язык C\# в качестве основного языка программирования, который значительно проще в освоении по сравнению с C++, применяемым в Unreal Engine, и при этом предоставляет достаточно высокую производительность . Среда разработки Unity отличается интуитивно понятным интерфейсом, а наличие огромного магазина активов с готовыми моделями, скриптами и инструментами позволяет существенно сократить время прототипирования . Эта экосистема делает Unity идеальным выбором для проектов, где важна скорость разработки и быстрая итеративность.

\subsubsection{Анализ возможностей Unreal Engine}
Unreal Engine, разрабатываемый компанией Epic Games, является одним из самых мощных и функционально насыщенных игровых движков в мире, занимающим лидирующие позиции в сегменте AAA-проектов, то есть игр с крупнейшими бюджетами и наиболее высокими требованиями к графике. В отличие от Unity, который часто выбирают за простоту и доступность, Unreal Engine позиционируется как инструмент для создания визуально впечатляющих проектов, где качество графики является первостепенной задачей. Анализ возможностей Unreal Engine важен для понимания того, какие архитектурные решения и технологии признаны индустрией как передовые, и в какой степени они могут быть адаптированы или осмыслены в рамках разработки собственного учебного движка.

Ключевым отличием Unreal Engine от большинства других движков является то, что его исходный код полностью открыт и доступен для изучения и модификации после подписания стандартного лицензионного соглашения. Это предоставляет разработчикам беспрецедентный уровень контроля: при необходимости можно изменить поведение любой подсистемы, от физики до рендеринга, и пересобрать движок под конкретные нужды проекта. Именно благодаря этой открытости Unreal Engine стал не только инструментом для создания игр, но и мощной образовательной платформой, на примере которой можно изучать устройство промышленного игрового движка. Для данной дипломной работы этот аспект важен, поскольку принципы, реализованные в Unreal Engine, во многом определяют ожидания от архитектуры игровых движков в целом.

В качестве основного языка программирования Unreal Engine использует C++, что является традиционным выбором для высокопроизводительных систем. Однако разработчики Epic Games внедрили мощную систему рефлексии, работающую через специальные макросы и генератор кода Unreal Header Tool, которая добавляет в C++ такие возможности, как сборка мусора, метаинформация о типах, сериализация и механизм событий. Благодаря этому разработчик получает производительность нативного кода с некоторыми удобствами, характерными для управляемых языков. Кроме того, Unreal Engine включает визуальный язык программирования Blueprints, позволяющий создавать игровую логику без написания кода, соединяя узлы, представляющие функции и переменные. Blueprints особенно полезны для дизайнеров уровней и художников, позволяя им быстро прототипировать поведение без участия программистов, а также служат отличным инструментом для визуализации потоков данных и логики.

\subsubsection{Анализ возможностей Godot}
Godot Engine представляет собой свободно распространяемый игровой движок с открытым исходным кодом, который за последние годы приобрёл значительную популярность, особенно среди независимых разработчиков и энтузиастов. В отличие от Unity и Unreal Engine, которые принадлежат коммерческим компаниям, Godot развивается силами сообщества и некоммерческого фонда Godot, что определяет его философию, лицензионную политику и архитектурные особенности. Анализ возможностей Godot важен для данной работы, поскольку этот движок демонстрирует альтернативный подход к разработке игровых движков, основанный на открытости, модульности и минимализме, и во многом близок по духу к учебному проекту, создаваемому в рамках диплома.

Ключевым преимуществом Godot является его лицензия MIT, которая позволяет использовать, модифицировать и распространять движок без каких-либо отчислений и без необходимости открывать собственный код производных проектов. Это делает Godot идеальным выбором для тех, кто хочет полностью контролировать свой технологический стек, не опасаясь изменений лицензионной политики, как это произошло с Unity в 2023 году. Для образовательных целей открытость Godot имеет особую ценность, поскольку студенты могут изучать исходный код движка, понимать внутреннее устройство его подсистем и даже вносить в него изменения. Этот аспект роднит Godot с целями данной дипломной работы, где самостоятельная разработка движка также направлена на глубокое понимание его внутренних механизмов.

Архитектурно Godot построен вокруг собственной системы узлов и сцен, которая существенно отличается от компонентной модели Unity или объектной модели Unreal Engine. В Godot всё является узлом, причём каждый узел выполняет конкретную функцию: узел KinematicBody отвечает за физическое тело, узел Sprite за отображение двухмерной графики, узел AudioStreamPlayer за воспроизведение звука. Сцена в Godot представляет собой дерево узлов, которое может быть сохранено в файл и затем использовано как подсистема внутри другой сцены. Это порождает иерархический подход к конструированию игровых объектов, когда сложное поведение собирается из простых узлов, вложенных друг в друга. Такой подход интуитивно понятен и позволяет быстро прототипировать структуру уровней, однако при глубокой вложенности может приводить к сложностям с отслеживанием связей между узлами.
